\section{Conclusiones del Capítulo}
\label{sec:cap4_conclusiones}

Este capítulo desarrolló la arquitectura híbrida que integra el modelo mecanicista DEB del Capítulo~\ref{chap:modelo_deb} con un calibrador neuronal paramétrico (MLP), resolviendo el problema de estados parcialmente observables que limitaba la aplicación operativa del modelo puro. A continuación se sintetizan los aportes principales.

\textbf{Calibrador neuronal paramétrico.} Se formuló y validó un MLP que opera como corrector de escala sobre el subconjunto de parámetros de primera sensibilidad $\{\tau_h, \rho_{\mathrm{conv}}, T_{\mathrm{opt}}\}$, sin recurrir a la predicción directa de biomasa. El MLP se entrena en esquema \emph{two-stage} (offline) con supervisión exclusiva del residuo de \ch{CO2} $\varepsilon_k = s_k - J_k^{\mathrm{DEB}}$, evitando el problema inverso mal condicionado y preservando la interpretabilidad fisiológica de cada parámetro. La salida acotada ($\pm 20\%$) garantiza que la red actúe como ajustador fino local, no como redefinidor de la biología larval.

\textbf{Estrategia de calibración por fases.} La calibración separó explícitamente los parámetros semilla fijos (Tabla~\ref{tab:parametros_semilla}) de los parámetros libres ajustados por L-BFGS-B con restricciones de caja (Tabla~\ref{tab:restricciones_caja}), culminando en la corrección neuronal en tiempo real mediante la ecuación~\eqref{eq:actualizacion_param}. Esta división preservó la interpretabilidad biofísica del núcleo mecanicista mientras adaptaba el modelo a las condiciones locales del reactor.

\textbf{Confrontación modelo--sensor y diagnóstico con \ch{CH4}.} El sistema híbrido predice la tasa de emisión de \ch{CO2} con MAPE $= 8{,}3\%$ y $R^2 = 0{,}78$, superando ampliamente el modelo lineal estático de referencia. El metano se mantuvo fuera de las ecuaciones de estado del DEB ---conforme al supuesto de aerobiosis estricta--- pero operó como canal paralelo de diagnóstico mediante un clasificador binario con umbral de \SI{100}{ppm}. Durante la validación experimental, el detector alcanzó un recall del $100\%$ en la identificación de eventos de anaerobiosis confirmados, validando la hipótesis~H4.

\textbf{Inventario dinámico de carbono.} La integración modelo--datos generó un inventario dinámico de carbono emitido por ciclo, cuantificando una discrepancia media acumulada $\bar{\Delta}_{\mathrm{Total}} = 5{,}1$~g~C ($3{,}6\%$ del total medido). Este residuo sistemático, atribuible a la respiración microbiana del sustrato no capturada por el DEB larval, es de magnitud reducida y permite su corrección mediante un factor de ajuste $\rho_{\mathrm{microb}} = 1{,}04$ en aplicaciones operativas futuras.

\textbf{Delimitación y puente.} Con la arquitectura híbrida validada y los parámetros calibrados, el siguiente capítulo traslada el sistema desde el dominio de la predicción hacia el dominio de la acción. Se desarrollará la plataforma de soporte a la decisión operativa, que integra: (i)~el semáforo de riesgo ambiental (verde/ámbar/rojo) alimentado por los residuos $\delta(t_k)$; (ii)~el protocolo de respuesta ante eventos de anaerobiosis; y (iii)~la generación automatizada de reportes MRV con trazabilidad de \emph{timestamp} y sellado de datos, orientados a la verificación externa del desempeño ambiental del proceso de bioconversión.

\subsection*{Cumplimiento del objetivo específico 3}

El objetivo específico 3 exigía desarrollar un algoritmo de estimación de estados mediante redes neuronales que infiriera la biomasa larval a partir de datos operativos. Este objetivo se cumple en tres dimensiones:

(i) \textit{Dimensión exploratoria.} Se implementó una red neuronal directa (RNA-B) que mapeaba observables $\mathbf{u}_k \mapsto \hat{B}_k$. La validación cruzada LORO ($k=6$ folds) reveló un fallo sistemático (MAPE validación $>40\%$, predicciones no físicas), demostrando que el problema inverso no es abordable mediante caja negra pura bajo muestreo destructivo limitado ($N\approx 30$). Se documenta en la Sección~4.6.

(ii) \textit{Dimensión de arquitectura alternativa.} Se evaluó conceptualmente una formulación UDE que integrara una red neuronal dentro de la dinámica continua del DEB. El análisis de identificabilidad, costo computacional y riesgo de sobreajuste inter-replica concluyó que la UDE no aporta ventaja predictiva ni interpretativa sobre el calibrador paramétrico validado, mientras introduce fragilidad numérica y complejidad desproporcionada. Esta vía fue descartada deliberadamente.

(iii) \textit{Dimensión operativa.} Se implementó un calibrador neuronal parámetrico (MLP) que corrige el subconjunto de primera sensibilidad $\{\tau_h, \rho_{\mathrm{conv}}, T_{\mathrm{opt}}\}$ del modelo DEB, utilizando como única supervisión el residuo de CO$_2$ $\varepsilon_k$. El MLP nunca recibe biomasa como entrada ni supervisión, delegando la dinámica biológica al modelo mecanicista. La biomasa estructural $B(t)$ se recupera indirectamente como salida del DEB re-simulado con parámetros corregidos. El sistema hiíbrido alcanza MAPE $= 8{,}3\%$ y $R^2 = 0{,}78$ en la predicción de emisiones, y MAPE$_B = 16{,}2\%$ en la recuperación de biomasa máxima en cosecha.

Por tanto, el objetivo específico 3 se considera cumplido: se desarrolló un algoritmo de estimación de estados mediante redes neuronales (el MLP calibrador), que infiere indirectamente el estado del sistema a través de la corrección paramétrica y mejora la precisión del modelo DEB para la estimación de CO$_2$. La trazabilidad completa del ajuste metodológico garantiza la robustez epistemológica de la solución adoptada.
